
% GitHub-Dateiname: MoE_Gloss_A_Normpunkt_261029-1.tex
\documentclass[a4paper,11pt]{article}

\usepackage[T1]{fontenc}
\usepackage[utf8]{inputenc}
\usepackage[ngerman]{babel}
\usepackage[scaled]{helvet}
\renewcommand{\familydefault}{\sfdefault}
\usepackage{sfmath}
\usepackage{amsmath}
\usepackage[a4paper,margin=25mm]{geometry}
\usepackage{microtype}

\setlength{\parindent}{0pt}
\setlength{\parskip}{0.8em}

\begin{document}

\begin{center}
{\Large\bfseries Apeiron-Normierungspunkt}

\vspace{0.4em}
MoE-Glossareintrag

Version 261029-1
\end{center}

\section*{Definition}

Der Apeiron-Normierungspunkt ist der Bezugspunkt
des normierten Apeiron-Expansionsprozesses.

An diesem Punkt gelten die Festlegungen

\[
\boxed{
\tau_A = V_A = \rho_A = A = 1
}
\]

Dabei bezeichnet $\tau_A$ den Prozessparameter,
$V_A$ das normierte Apeiron-Volumen,
$\rho_A$ die normierte Apeiron-Dichte
und $A$ die erhaltene Apeiron-Menge.

\section*{Mathematische Festlegung}

Für den Apeiron-Prozess gilt die Erhaltung

\[
A = \rho_A V_A = \mathrm{konstant}.
\]

Mit der Normierung $A=1$ folgt

\[
\rho_A = \frac{1}{V_A}.
\]

Der Prozessparameter wird proportional
zum normierten Apeiron-Volumen festgelegt:

\[
\boxed{
\tau_A = V_A = \frac{1}{\rho_A}
}
\]

Damit beschreibt $\tau_A$ den Fortschritt
des Expansionsprozesses anhand des
normierten Apeiron-Volumens.

\section*{Prozessgrenzen}

Für den Anfangsgrenzzustand gilt

\[
\tau_A \to 0,\qquad
V_A \to 0,\qquad
\rho_A \to \infty.
\]

Für den Endgrenzzustand gilt

\[
\tau_A \to \infty,\qquad
V_A \to \infty,\qquad
\rho_A \to 0.
\]

Der Normierungspunkt liegt bei

\[
(\tau_A,V_A,\rho_A)=(1,1,1).
\]

\section*{Verweise}

Apeiron; Apeiron-Expansion; Prozessparameter.

\vfill
\begin{center}
\small
Formulierung, Recherche und Prüfung KI-unterstützt
\end{center}

\end{document}
